\section{The Admission Grammar}
\label{sec:grammar}

The grammar names what a substrate is willing to admit. We describe it informally and cite the underlying formalisation in companion work~\cite{programmableSovereignty2026}. The point is to give the reader a clear picture of the typed object that admission operates on; the load-bearing theorems live in the cited substrate.

\paragraph{Action classes.}
A tool call belongs to one of three action classes. \emph{Observational} calls read state without mutating it (a database SELECT, a file read, an idempotent HTTP GET). \emph{Reversible} calls mutate state with a known inverse executor: a file quarantine has a release, a row insert has a delete, a process suspend has a resume. \emph{Destructive} calls mutate state without a typed inverse: a permanent record deletion that empties a backup, a transaction broadcast to a public ledger, an email sent to an external recipient.

The classification is action-dependent rather than agent-dependent, and is the substrate's commitment, not the agent's. The substrate decides, for each tool a server exposes, which class each method falls into. An unclassified method is destructive by default (fail-closed).

\paragraph{The admission envelope.}
The substrate accepts a tool call wrapped in an admission envelope that carries four fields beyond the call's intrinsic arguments: (i) an authorisation chain (the capability token under which the agent acts), (ii) an action class (observational, reversible, or destructive), (iii) a positive TTL by construction (the maximum duration the call's effects may stand before the substrate considers them stale), and (iv) for reversible and destructive classes, a rollback receipt slot (an opaque record reserved at admission to be filled at rollback time).

The envelope is typed. A reversible call without a TTL is not constructable: the substrate's input parser rejects it before any executor sees it. A destructive call without both a TTL and a bilateral cosignature is not constructable. The envelope is the syntactic constraint, not a runtime policy.

\paragraph{The four operations.}
Admission proceeds by four decidable operations. (1) \emph{Capability check}: the authorisation chain attenuates a valid trust root and the residual permits the call's intended effect on the named subject. (2) \emph{Class check}: the call's classified action class matches the envelope's declared class. (3) \emph{TTL check}: the envelope's declared TTL is positive and bounded by the substrate's per-class maximum. (4) \emph{Rollback slot}: for reversible and destructive classes, the envelope carries a typed rollback receipt slot; for destructive classes, a bilateral cosignature is also present from both the agent's operator polity and the tool server's host polity.

\paragraph{Structural theorems carried by the substrate.}
The grammar carries three structural commitments discharged in companion work~\cite{programmableSovereignty2026}.

\thm{bounded\_executive\_action\_carries\_ttl\_and\_rollback\_slot}: a reversible or destructive call's envelope is well-typed only if it carries a positive TTL witness and a rollback receipt slot. The theorem is a definitional bridge documenting the type-level invariant; its substantive use is the type-checker's refusal to construct envelopes without the discipline.

\thm{rollback\_admission\_composes\_with\_refinement}: when the admission predicate is amended (the policy changes), a rollback receipt that closes a previously-admitted call remains admissible under the amended predicate, provided the amendment is backward-refining. Rollback obligations do not vanish across policy changes.

\thm{treaty\_admission\_iff\_predicate\_intersection}: a destructive call admits only if both the operator polity and the host polity admit it under their respective predicates. This is the structural reason a single-party operator manipulation does not suffice to admit a destructive tool call against a counterparty whose policy forbids it.

\paragraph{What the grammar does not claim.}
The grammar does not claim that every tool call should be destructive-class. Observational and reversible classes are the common case and admit on much weaker grounds; the discipline matters at the destructive boundary. The grammar does not claim that the substrate can synthesise a rollback executor for an intrinsically irreversible call. It claims only that, for such a call, admission requires bilateral cosignature in addition to the typed witness, and the cosignature is a public commitment that the call has been admitted as irreversible by both parties.

The central commitment is that a destructive tool call is not constructable in the typed runtime path without a rollback witness or its destructive-equivalent (bilateral cosignature). Misalignment in the issuing model does not produce envelopes that bypass the host polity's independent admission. The two-party structural check, not the operator's diligence, is what the substrate enforces.
